\appendix
\section*{Appendix A. Regional update constants}
For search dimension $n$ and effective parent count $m=\mu_{\mathrm{eff}}$, the learning rates and damping factor are
\Eq{cc}{c_c=\frac{4+m/n}{n+4+2m/n}.}
\Eq{cs}{c_\sigma=\frac{m+2}{n+m+5}.}
\Eq{c1}{c_1=\frac{2}{(n+1.3)^2+m}.}
\Eq{cmu}{c_\mu=\min\left(1-c_1,\,\frac{2(m-2+1/m)}{(n+2)^2+m}\right).}
\Eq{damp}{d_\sigma=1+2\max\left(0,\sqrt{\frac{m-1}{n+1}}-1\right)+c_\sigma.}
\Eq{chi}{\chi_n=\sqrt n\left(1-\frac{1}{4n}+\frac{1}{21n^2}\right).}
Here $\chi_n$ approximates the expected norm of a standard Gaussian vector. Eigenvalues are floored at $10^{-14}$ when forming the square root and inverse square root of the covariance. The upper end of the selector coordinate in the lower region is clipped to $1-10^{-12}$ before affine mapping, to keep boundary points in their assigned region. If the last batch contains fewer than $\mu$ trials, its evaluations have already entered the archives and the distribution update is omitted.

\section*{Appendix B. Convergence of the best admissible cost}
Retain the parameter regions, numerical feasible sets and represented optimum of Section~3. Consider an infinite run with a deterministic decoder and rank. Let $\nu_b$ denote the law of a fresh local uniform draw mapped to region $b$, including the boundary handling of Appendix A. Assume that the reserved draws are mutually independent. Let $C_B$ denote the best admissible cost after $B$ evaluations, set to $+\infty$ before the first admission. For $\epsilon>0$, define the near-optimal subset and its probability by
\Eq{near}{A_{b,\epsilon}=\{\xi\in\mathcal F_b:J(D(\xi))<J_D^\star+\epsilon\}.}
\Eq{mass}{v_{b,\epsilon}=\nu_b(A_{b,\epsilon}).}

\textbf{Theorem 1.} Let $J_D^\star$ be finite. Suppose that for each $\epsilon>0$ at least one region $b\in\mathcal I$ has $v_{b,\epsilon}>0$. With one independent refresh in every reserved batch, the best admissible cost converges almost surely to $J_D^\star$. At budget $B=N_0+\lambda N$, after $N$ full batches, its upper-tail probability satisfies
\Eq{bound}{\Pr\{C_B\ge J_D^\star+\epsilon\}\le\prod_{b\in\mathcal I}(1-v_{b,\epsilon})^{\lfloor N/(HK)\rfloor}.}
A factor with zero exponent is interpreted as one.

\textit{Proof.} Each full block of $HK$ batches contains exactly $K$ reservations. Cyclic assignment sends one of these batches to each region. Consequently each region receives at least $\lfloor N/(HK)\rfloor$ independent refreshes among the first $N$ batches. A draw in $A_{b,\epsilon}$ has zero violation and cost below $J_D^\star+\epsilon$. The lexicographic acceptance rule retains that cost or a smaller one for the rest of the run. The event on the left of Eq.~(\ref{eq:bound}) therefore entails that every reserved draw missed its corresponding set. Independence gives the product bound. This step requires independence only for the reserved draws.

For each integer $j\ge1$, take $\epsilon=1/j$. By hypothesis the right-hand side of Eq.~(\ref{eq:bound}) tends to zero. The event that the archive never reaches this neighborhood is contained in every finite-budget miss event and hence has probability zero. Taking the countable intersection of the probability-one hit events shows that, almost surely, all these neighborhoods are eventually reached. The archive costs are nonincreasing and bounded below by $J_D^\star$ once finite. They therefore converge to $J_D^\star$. \hfill$\square$

The positive-mass assumption connects the proof to the engineering problem. For example, an admissible parameter with a feasible interior neighborhood and a continuous objective admits positive-probability draws arbitrarily close to its cost if the refresh law has positive density there. A terminal equality applied solely after forward integration can instead restrict admissible parameters to a set of zero volume. A maneuver omitted by the decoder cannot contribute to $J_D^\star$. The theorem concerns the represented admissible set and an ideal independent-draw process. A finite pseudorandom run and its numerical replay support separate empirical claims about a stated budget and acceptance tests.

